\documentclass[12pt]{article}

\usepackage[paperwidth=612pt,paperheight=792pt,top=72pt,bottom=72pt,left=30mm,right=20mm]{geometry}
\usepackage{times}
\usepackage{setspace}
\usepackage{fancyhdr}
\usepackage{ragged2e}
\usepackage{enumitem}
\setstretch{1.5}

\pagestyle{fancy}
\fancyhf{}
\fancyhead[R]{AyurSpace}
\fancyfoot[L]{Pillai HOC College of Engineering \& Technology (Autonomous), B. E. Computer Engineering}
\fancyfoot[R]{\thepage}
\renewcommand{\headrulewidth}{1pt}
\renewcommand{\footrulewidth}{1pt}

\begin{document}

\vspace*{\fill}
\begin{center}
\textbf{\Huge{Chapter 9}}\\
\vspace{5mm}
\textbf{\Huge{Conclusion \& Future Work}}
\end{center}
\vspace*{\fill}
\newpage

\section*{9.1 \quad Conclusion}
\addcontentsline{toc}{section}{9.1 Conclusion}
\justify

The creation of \textbf{AyurSpace} is a key step in the digital preservation and democratization
of Ayurvedic knowledge. By effectively executing a hybrid neuro-symbolic architecture, this
framework deals with the major problems facing current plant identification tools---specifically,
the absence of medicinal context and the risk of generative hallucinations. Our experimental
findings, producing a 96.4\% accuracy in taxonomy and a 99.1\% compliance in contextual generative
reasoning, confirm the effectiveness of decoupling visual recognition from semantic analysis. This
dependability is extremely important in the health field, where misidentification can have
toxicological effects. Additionally, the algorithmic formalization of the \textit{Dosha} assessment
transforms subjective clinical intuition into a replicable digital logic, enabling users to make
smart health decisions based on ancient wisdom merged with robust modern mobile processing.

\section*{9.2 \quad Limitations}
\addcontentsline{toc}{section}{9.2 Limitations}
\justify

\begin{enumerate}
    \item \textbf{Network Dependency}: The current architecture requires an active internet
    connection for both Plant.id and Gemini API calls. Users in areas with poor connectivity
    experience degraded performance or complete failure of AI features.

    \item \textbf{API Cost}: Both Plant.id and Google Gemini APIs operate on usage-based pricing
    models. High-volume production usage would incur significant operational costs.

    \item \textbf{Dataset Bias}: Plant.id's training dataset may have lower representation for
    rare Indian medicinal herbs, potentially reducing accuracy for uncommon species.

    \item \textbf{LLM Hallucinations}: While significantly mitigated by the hybrid pipeline,
    Gemini may occasionally produce inaccurate Ayurvedic properties for plants with limited
    online documentation.
\end{enumerate}

\newpage

\section*{9.3 \quad Future Scope}
\addcontentsline{toc}{section}{9.3 Future Scope}
\justify

Future research and scale expansions will focus on three main development branches:

\begin{enumerate}
    \item \textbf{Edge AI Optimization}: Utilizing quantized MobileNet models (int8 quantization)
    on the mobile device natively. This will allow for the complete offline visual identification
    of plants, resolving accessibility issues for users in remote rural locations suffering from
    weak network connectivity.

    \item \textbf{Augmented Reality (AR) Interfaces}: Harnessing ARCore frameworks to overlay
    medicinal properties, alerts, and detailed plant structures directly onto the camera
    viewfinder. This transition to augmented realities will cultivate deeply immersive, real-time
    educational experiences for students and the general public.

    \item \textbf{Telemedicine Integration}: Launching a ``Vaidya Connect'' application component.
    By piping users' generated Dosha baselines and scan histories into professional Ayurvedic
    dashboards, we can create an end-to-end telemedicine ecosystem that blends self-evaluation
    algorithms with direct oversight from licensed Ayurvedic doctors.
\end{enumerate}

\end{document}
